\documentclass[11pt,a4paper]{article}
\usepackage[margin=2.2cm,top=2.0cm,bottom=2.2cm]{geometry}
\usepackage{amsmath,amssymb}
\usepackage[T1]{fontenc}
\usepackage{mathptmx}
\usepackage{xcolor}
\usepackage{enumitem}
\usepackage{fancyhdr}
\usepackage{tabularx}
\usepackage{mdframed}
\usepackage{multicol}

\definecolor{navy}{HTML}{1F4E79}
\definecolor{rule}{HTML}{333333}
\definecolor{boxbg}{HTML}{F4F6F9}
\definecolor{keyred}{HTML}{B02020}

% \ANSWERStrue for the tutor key
\newif\ifANSWERS
\input{mode}

\pagestyle{fancy}\fancyhf{}
\renewcommand{\headrulewidth}{0pt}
\renewcommand{\footrulewidth}{0.4pt}
\fancyfoot[L]{\footnotesize\color{navy}SAT Math: Foundations to Mastery \ $\cdot$\ Systems and Inequalities\ifANSWERS\ $\cdot$\ \textbf{\color{keyred}ANSWER KEY}\fi}
\fancyfoot[R]{\footnotesize\color{navy}Page \thepage\ of \pageref{LastPage}}
\usepackage{lastpage}

\setlength{\parindent}{0pt}
\setlength{\parskip}{4pt}

\newcommand{\PartHead}[2]{%
  \vspace{6pt}
  {\color{navy}\rule{\textwidth}{0.8pt}}\\[2pt]
  {\large\bfseries\color{navy}#1}\hfill{\small\itshape #2}\\[-4pt]
  {\color{navy}\rule{\textwidth}{0.4pt}}
  \vspace{4pt}
}
\newcommand{\sitting}[1]{{\footnotesize\textsc{\color{navy}#1}}\par\vspace{2pt}}
\newcommand{\spr}{\\[3pt]\hspace*{6pt}Answer: \underline{\hspace{3.2cm}}\hfill{\footnotesize\itshape(write your answer)}}
\ifANSWERS
  \newcommand{\ans}[1]{\ {\color{keyred}\bfseries #1}}
  \newcommand{\note}[1]{\par{\footnotesize\color{keyred}#1}}
\else
  \newcommand{\ans}[1]{}
  \newcommand{\note}[1]{}
\fi

\newcommand{\method}{\\[3pt]\hspace*{6pt}Method: \underline{\hspace{4.6cm}}\hfill{\footnotesize\itshape(substitution, elimination, or graphing)}}

\begin{document}
\label{firstpage}

\begin{center}
{\LARGE\bfseries\color{navy}Systems and Inequalities}\\[4pt]
{\color{navy}SAT Math $\cdot$ Foundations to Mastery $\cdot$ Week of 11 October}
\ifANSWERS\\[4pt]{\color{keyred}\bfseries ANSWER KEY --- tutor copy}\fi
\end{center}
\vspace{2pt}
{\color{rule}\rule{\textwidth}{1pt}}
\vspace{6pt}

\ifANSWERS\else
Name: \underline{\hspace{6cm}} \hfill Due: \underline{\hspace{3.2cm}}
\vspace{6pt}
\fi

\begin{mdframed}[backgroundcolor=boxbg,linecolor=navy,linewidth=0.6pt,innertopmargin=6pt,innerbottommargin=8pt]
{\bfseries\color{navy}How to use this sheet}\par\vspace{2pt}
This is four short sittings, not one long one. Four half-hours across the week will
do more for you than one long Saturday, and you are far more likely to finish it.

\begin{itemize}[leftmargin=14pt,itemsep=1pt,topsep=3pt]
  \item \textbf{Part A is no calculator.} None. It is 15 minutes and it keeps the last two weeks alive.
  \item \textbf{In Part B, name your method before you start.} Every time. Choosing the cheapest route is the skill being graded here, not grinding.
  \item When you get one wrong, write it in your error log: what you did, and what you should have done.
\end{itemize}
\end{mdframed}

\vspace{4pt}
\begin{center}
\begin{tabularx}{\textwidth}{@{}l l X r@{}}
\textbf{Sitting} & \textbf{Part} & \textbf{What it drills} & \textbf{Minutes}\\[2pt]
\hline\rule{0pt}{11pt}
Sunday   & A    & Retention: signs and slope                       & 15\\
Tuesday  & B    & Solve the system, and name the method            & 25\\
Thursday & C    & Inequalities, including translation              & 25\\
Saturday & D, E & Build the system from words; six mixed, timed    & 25\\
\hline
\end{tabularx}
\end{center}

% ------------------------------------------------ PART A
\PartHead{Part A --- Retention}{No calculator. 15 minutes.}
\sitting{Sunday}
The first five keep last fortnight's sign work alive. The last five keep week two's
lines alive. Short answers. If a sign is in doubt, draw a number line in the margin.

\begin{enumerate}[leftmargin=20pt,itemsep=7pt]
\item $8-(-5)=$ \ifANSWERS\else\underline{\hspace{2cm}}\fi \ans{13}
\item $(-4)^2-4^2=$ \ifANSWERS\else\underline{\hspace{2cm}}\fi \ans{0 \ (both squares are $+16$)}
\item Expand: $-2(x-6)=$ \ifANSWERS\else\underline{\hspace{4cm}}\fi \ans{$-2x+12$}
\item $f(x)=x^2-3x$. \ Find $f(-2)$: \ifANSWERS\else\underline{\hspace{2cm}}\fi \ans{$(-2)^2-3(-2)=4+6=10$}
\item $-(5-9)=$ \ifANSWERS\else\underline{\hspace{2cm}}\fi \ans{4}
\item A line passes through $(-2,7)$ and $(4,-5)$. Its slope is \ifANSWERS\else\underline{\hspace{2cm}}\fi \ans{$\dfrac{-5-7}{4-(-2)}=\dfrac{-12}{6}=-2$}
\item For $3x-2y=12$, the $y$-coordinate of the $y$-intercept is \ifANSWERS\else\underline{\hspace{2cm}}\fi \ans{$-6$ \ (put $x=0$: $-2y=12$)}
\item For $3x-2y=12$, the slope is \ifANSWERS\else\underline{\hspace{2cm}}\fi \ans{$1.5$ \ ($y=1.5x-6$)}
\item For $y=4-0.5x$, the slope is \ifANSWERS\else\underline{\hspace{2cm}}\fi \ans{$-0.5$}
\item A repair costs $C(t)=45+18t$ dollars for $t$ hours of labour. In one short phrase, what does the 18 represent?\ifANSWERS\else\\
  \underline{\hspace{\linewidth}}\fi
  \ans{\\ The cost per hour of labour --- the rate the bill grows for each extra hour.}
\end{enumerate}

\note{\textbf{Marking rule: below 8 of 10 means an earlier week has decayed.} Items 2, 5 and 6 are the diagnostic ones for signs; items 7 and 8 separate a student who reads standard form from one who guesses at it.}

\ifANSWERS\else\newpage\fi
% ------------------------------------------------ PART B
\PartHead{Part B --- Solve the system, and name the method}{Calculator allowed. 25 minutes.}
\sitting{Tuesday}
You have three tools: substitution, elimination, and graphing both lines on Desmos.
For each one, write the answer \emph{and} write which tool you chose. Choosing is a
decision, and on a timed test the cheapest tool is worth as much as the arithmetic.
Look at the system for five seconds before you pick up your pen.

\begin{enumerate}[leftmargin=20pt,itemsep=9pt]
\item $y=2x-3$ \ and \ $4x+y=15$. \ What is the value of $x$? \spr\method
  \ans{\\ $x=3$ (and $y=3$). Substitution is cheapest: $y$ is already alone, so $4x+(2x-3)=15$, $6x=18$.}

\item $x=4-3y$ \ and \ $2x+5y=3$. \ What is the value of $y$? \spr\method
  \ans{\\ $y=5$ (and $x=-11$). Substitution is cheapest: $x$ is already alone, so $2(4-3y)+5y=3$, $8-y=3$.}

\item $4x+3y=18$ \ and \ $4x-5y=2$. \ What is the value of $y$? \spr\method
  \ans{\\ $y=2$ (and $x=3$). Elimination is cheapest: the $4x$ terms match, so subtracting gives $8y=16$ in one line.}

\item $3x+2y=7$ \ and \ $5x-2y=17$. \ What is the value of $y$? \spr\method
  \ans{\\ $y=-1$ (and $x=3$). Elimination is cheapest: $+2y$ and $-2y$ are opposites, so adding gives $8x=24$. Then $2y=7-9=-2$.}

\item $0.6x+1.5y=6.9$ \ and \ $2.4x-0.5y=8.1$. \ What is the value of $x$? \spr\method
  \ans{\\ $x=4$ (and $y=3$). Graphing is cheapest: the decimals make either algebraic route slow, and Desmos reads off the intersection $(4,3)$ with no risk of a decimal slip.}

\item $7x+11y=101$ \ and \ $11x+7y=97$. \ What is the value of $x+y$? \spr\method
  \ans{\\ $x+y=11$. Add the two equations as they stand: $18x+18y=198$, so $x+y=11$. You never need $x$ or $y$ separately. (Had you been asked to solve it, graphing would have been cheapest, because clearing 7s and 11s by hand is slow.)}
  \note{\textbf{This is the highest-value item on the sheet.} The question asks for a combination, not for the variables, and the combination falls straight out of adding the equations. A student who solves the whole system still gets 11, but burns two minutes to do it. If he spotted this, say so out loud. For the record, $x=5$ and $y=6$.}

\item How many solutions does the system $2x-6y=9$ and $x-3y=7$ have?\\
  \makebox[3.9cm][l]{A) Zero}\makebox[3.9cm][l]{B) Exactly one}%
  \makebox[3.9cm][l]{C) Exactly two}\makebox[3.9cm][l]{D) Infinitely many}
  \method
  \ans{\\ A. Doubling the second equation gives $2x-6y=14$, but the first says $2x-6y=9$. Same left side, different right side, so the lines are parallel and never meet. Graphing confirms it in seconds. \\ B is the assumption that every pair of lines crosses. D comes from noticing the coefficients are proportional and stopping before checking the constants. C is a quadratic habit carried into a linear question.}

\item How many solutions does the system $3x+4y=10$ and $9x+12y=30$ have?\\
  \makebox[3.9cm][l]{A) Zero}\makebox[3.9cm][l]{B) Exactly one}%
  \makebox[3.9cm][l]{C) Exactly two}\makebox[3.9cm][l]{D) Infinitely many}
  \method
  \ans{\\ D. Multiplying the first equation by 3 gives the second exactly, so it is one line written twice and every point on it works. \\ A comes from seeing the proportional coefficients and calling the lines parallel without checking the constants --- the exact mirror of item 7. B is the default assumption again.}

\end{enumerate}

\ifANSWERS\else\newpage\fi
% ------------------------------------------------ PART C
\PartHead{Part C --- Inequalities}{Calculator allowed. 25 minutes.}
\sitting{Thursday}
Two things are being tested: the flip when you multiply or divide by a negative, and
reading an English phrase into a symbol. After every solve, test one value on each
side of the boundary. That check costs ten seconds and catches a missing flip every time.

\begin{enumerate}[leftmargin=20pt,itemsep=9pt]
\item Which inequality gives all values of $x$ for which $5-3x\ge17$?\\
  \makebox[3.6cm][l]{A) $x\ge-4$}\makebox[3.6cm][l]{B) $x\le-4$}%
  \makebox[3.6cm][l]{C) $x\le4$}\makebox[3.6cm][l]{D) $x\ge4$}
  \ans{\\ B. $-3x\ge12$, then divide by $-3$ \emph{and flip}. Check: $x=-5$ gives $20\ge17$, true; $x=-3$ gives $14\ge17$, false. \\ A is the forgotten flip. C drops the minus when dividing. D does both.}

\item What is the greatest integer $x$ for which $4-7x\ge-10$? \spr
  \ans{\\ $2$. $-7x\ge-14$, divide by $-7$ and flip: $x\le2$. Check: $x=2$ gives $-10\ge-10$, true; $x=3$ gives $-17\ge-10$, false.}

\item Which inequality gives all values of $x$ for which $-2(x-5)>14$?\\
  \makebox[3.9cm][l]{A) $x<-2$}\makebox[3.9cm][l]{B) $x>-2$}%
  \makebox[3.9cm][l]{C) $x<-12$}\makebox[3.9cm][l]{D) $x>-12$}
  \ans{\\ A. $-2x+10>14$, so $-2x>4$, then flip: $x<-2$. Check: $x=-3$ gives $16>14$, true; $x=-1$ gives $12>14$, false. \\ B is the forgotten flip. C comes from distributing as $-2x-10$, which gives $-2x>24$. D makes both errors.}

\item A game gives you 4 points for each task you finish. You already have 14 points and you need a total of at least 50. Which inequality describes the number of tasks $t$ that is enough?\\
  \makebox[3.9cm][l]{A) $14+4t>50$}\makebox[3.9cm][l]{B) $14+4t\ge50$}%
  \makebox[3.9cm][l]{C) $14+4t\le50$}\makebox[3.9cm][l]{D) $4t\ge50$}
  \ans{\\ B. ``At least 50'' means 50 counts, so the symbol is $\ge$. (It gives $t\ge9$.) \\ A reads ``at least'' as ``more than'' and loses the boundary. C reverses the direction. D forgets the 14 points already banked.}

\item A food truck charges \$3 for a taco and \$5 for a burrito. On one shift it took in no more than \$150 from these two items. Which combination is possible?\\
  \makebox[5.7cm][l]{A) 20 tacos and 20 burritos}\makebox[5.7cm][l]{B) 30 tacos and 12 burritos}\\[2pt]
  \makebox[5.7cm][l]{C) 40 tacos and 8 burritos}\makebox[5.7cm][l]{D) 10 tacos and 25 burritos}
  \ans{\\ B. The condition is $3t+5b\le150$. B gives $90+60=150$, and ``no more than'' includes 150. \\ A gives 160, C gives 160, D gives 155 --- all over. A student who reads ``no more than'' as ``less than'' rejects B as well and is left with nothing, which is the signal that the boundary was mishandled.}

\item Five less than the product of $-4$ and a number is more than 23. Which inequality describes all such numbers $n$?\\
  \makebox[3.9cm][l]{A) $n<-7$}\makebox[3.9cm][l]{B) $n>-7$}%
  \makebox[3.9cm][l]{C) $n<-4.5$}\makebox[3.9cm][l]{D) $n>-4.5$}
  \ans{\\ A. The sentence is $-4n-5>23$, so $-4n>28$, then flip: $n<-7$. Check: $n=-8$ gives $27>23$, true; $n=-6$ gives $19>23$, false. \\ B is the forgotten flip. C and D come from subtracting the 5 instead of adding it, giving $-4n>18$.}

\item On every day of a week the high temperature stayed between 12 and 19 degrees, and it never reached either of those two values. Which inequality describes the high temperature $t$?\\
  \makebox[3.9cm][l]{A) $12\le t\le19$}\makebox[3.9cm][l]{B) $12<t<19$}%
  \makebox[3.9cm][l]{C) $12>t>19$}\makebox[3.9cm][l]{D) $t>12$ or $t<19$}
  \ans{\\ B. ``Never reached either value'' makes both ends strict. \\ A includes the endpoints the sentence rules out. C writes the chain backwards, which describes nothing at all. D uses ``or'', which lets in every number.}

\item A line is drawn on the $xy$-plane through $(0,-3)$ with slope 2. The line itself is dashed, and the whole region \emph{above} it is shaded. Which inequality does the shaded region represent?\\
  \makebox[3.9cm][l]{A) $y>2x-3$}\makebox[3.9cm][l]{B) $y\ge2x-3$}%
  \makebox[3.9cm][l]{C) $y<2x-3$}\makebox[3.9cm][l]{D) $y\le2x-3$}
  \ans{\\ A. The boundary is $y=2x-3$; above means $y$ is larger, and dashed means the line is excluded. Test $(0,0)$: $0>-3$, true, and $(0,0)$ does lie above the line. \\ B ignores the dashed line. C picks the wrong side. D makes both errors.}

\end{enumerate}

\ifANSWERS\else\newpage\fi
% ------------------------------------------------ PART D
\PartHead{Part D --- Build the system from words}{No calculator. 15 minutes.}
\sitting{Saturday}
\textbf{Write the two equations only. Do not solve them.} That is the whole task.
Setting up the pair is the step that costs you marks, so this part isolates it and
nothing else. Say in words what each letter stands for, then write the two equations.

\begin{enumerate}[leftmargin=20pt,itemsep=11pt]
\item You buy 14 items in total, some pens at \$2 each and some notebooks at \$5 each, and you spend \$43 altogether. Let $p$ be the number of pens and $n$ the number of notebooks.\\
  \underline{\hspace{\linewidth}}\\[4pt]
  \underline{\hspace{\linewidth}}
  \ans{\\ $p+n=14$ \ (a count equation) \ and \ $2p+5n=43$ \ (a money equation). One equation counts things, the other counts dollars --- look for that pair every time.}

\item Two numbers have a sum of 56. One of them is 8 more than twice the other. Let $x$ be the larger number and $y$ the smaller.\\
  \underline{\hspace{\linewidth}}\\[4pt]
  \underline{\hspace{\linewidth}}
  \ans{\\ $x+y=56$ \ and \ $x=2y+8$. ``8 more than twice $y$'' is $2y+8$, not $2(y+8)$.}

\item A school sold 160 tickets to a concert. Adult tickets were \$12 and child tickets were \$7, and the takings were \$1{,}395. Let $a$ be the number of adult tickets and $c$ the number of child tickets.\\
  \underline{\hspace{\linewidth}}\\[4pt]
  \underline{\hspace{\linewidth}}
  \ans{\\ $a+c=160$ \ and \ $12a+7c=1395$. Again: one count equation, one money equation.}

\item A rectangular garden has a perimeter of 86 metres. Its length is 7 metres more than twice its width. Let $L$ be the length and $W$ the width, both in metres.\\
  \underline{\hspace{\linewidth}}\\[4pt]
  \underline{\hspace{\linewidth}}
  \ans{\\ $2L+2W=86$ \ and \ $L=2W+7$. Writing $L+W=86$ for the perimeter is the usual slip here.}

\item Maya has 18 coins, all dimes and quarters, worth \$3.15 in total. Let $d$ be the number of dimes and $q$ the number of quarters.\\
  \underline{\hspace{\linewidth}}\\[4pt]
  \underline{\hspace{\linewidth}}
  \ans{\\ $d+q=18$ \ and \ $0.10d+0.25q=3.15$. Mixing units is the trap: either work in dollars throughout, as here, or in cents throughout with $10d+25q=315$. Never both.}

\end{enumerate}

\ifANSWERS\else\newpage\fi
% ------------------------------------------------ PART E
\PartHead{Part E --- Mixed, timed}{Calculator allowed. 12 minutes, timed.}
\sitting{Saturday}
Set a timer for 12 minutes. These span Parts B to D. Answer every one --- there is no
penalty for a wrong answer, so a blank is a mark thrown away. Then check each answer
by substituting it back into both equations.

\begin{enumerate}[leftmargin=20pt,itemsep=9pt]
\item $y=3x-8$ \ and \ $2x+y=7$. \ What is the value of $y$? \spr
  \ans{\\ $y=1$. Substitution: $2x+(3x-8)=7$, so $5x=15$ and $x=3$, giving $y=3(3)-8=1$. Check in the second equation: $6+1=7$.}

\item $6x+5y=37$ \ and \ $5x+6y=29$. \ What is the value of $x+y$? \spr
  \ans{\\ $6$. Add the equations: $11x+11y=66$, so $x+y=6$. No need to find $x$ and $y$ separately (they are 7 and $-1$).}

\item Which inequality gives all values of $x$ for which $9-4x\le29$?\\
  \makebox[3.6cm][l]{A) $x\ge-5$}\makebox[3.6cm][l]{B) $x\le-5$}%
  \makebox[3.6cm][l]{C) $x\ge5$}\makebox[3.6cm][l]{D) $x\le5$}
  \ans{\\ A. $-4x\le20$, divide by $-4$ and flip. Check: $x=0$ gives $9\le29$, true; $x=-6$ gives $33\le29$, false. \\ B is the forgotten flip. C drops the minus when dividing. D does both.}

\item A theatre sold 90 tickets in total. Adult tickets cost \$15 and student tickets cost \$9, and the takings were \$1{,}062. How many adult tickets were sold? \spr
  \ans{\\ $42$. From $a+s=90$ and $15a+9s=1062$: substituting $s=90-a$ gives $15a+810-9a=1062$, so $6a=252$. Then $s=48$, and $15(42)+9(48)=630+432=1062$.}

\item A tutor charges a \$40 booking fee plus \$28 for each hour. A family wants the total to be no more than \$180. Which inequality describes the number of hours $h$ they can book?\\
  \makebox[3.9cm][l]{A) $40+28h<180$}\makebox[3.9cm][l]{B) $40+28h\le180$}%
  \makebox[3.9cm][l]{C) $28h\le180$}\makebox[3.9cm][l]{D) $40+28h\ge180$}
  \ans{\\ B. ``No more than \$180'' includes \$180, so the symbol is $\le$. (It gives $h\le5$.) \\ A loses the boundary. C forgets the booking fee. D reverses the direction.}

\item In the system $2x+ky=10$ and $6x+9y=30$, the constant $k$ is chosen so that the system has infinitely many solutions. What is the value of $k$? \spr
  \ans{\\ $3$. Multiplying the first equation by 3 gives $6x+3ky=30$, which must match $6x+9y=30$ term for term, so $3k=9$. With $k=3$ the two equations are the same line.}

\end{enumerate}

\ifANSWERS\else
\vspace{10pt}
\PartHead{Error log}{Fill this in as you go.}
\begin{center}
\begin{tabularx}{\textwidth}{|c|c|X|X|}
\hline
\textbf{Date} & \textbf{Q} & \textbf{What I did} & \textbf{What I should have done}\\
\hline
\rule{0pt}{22pt} & & &\\ \hline
\rule{0pt}{22pt} & & &\\ \hline
\rule{0pt}{22pt} & & &\\ \hline
\rule{0pt}{22pt} & & &\\ \hline
\rule{0pt}{22pt} & & &\\ \hline
\rule{0pt}{22pt} & & &\\ \hline
\rule{0pt}{22pt} & & &\\ \hline
\end{tabularx}
\end{center}
\vspace{4pt}
\begin{center}\small\itshape
Bring this sheet and your error log to the next session.
\end{center}
\fi

\ifANSWERS
\vspace{8pt}
\PartHead{Five-minute marking routine}{Tutor}
\begin{enumerate}[leftmargin=20pt,itemsep=4pt]
\item \textbf{Score Part A only.} Below 8/10 means an earlier week has decayed, and the next session's warm-up repairs that week rather than moving on.
\item \textbf{In Part B, check two things: whether he named a method at all, and whether he picked the cheapest one.} Grinding through substitution on every item is a miss even when all eight answers are right, because on a timed test it costs him the minutes he needs at the end. Item 6 is the one to ask about first.
\item \textbf{Pick two wrong answers from Part E} and have him narrate his thinking out loud. Leave the rest unmarked.
\end{enumerate}
If the error log is filled in, say so out loud. He will do more of what you visibly value.
\fi

\end{document}
